\section{Agentic Workflow vs. Autonomous Agent}

\subsection{概念区分}

\begin{knowledgebox}{Anthropic 对 Agent 的定义}
根据 Anthropic 年初关于 Agent 的定义，Plan-Execute-RePlan 属于\textbf{预定义的 Agentic Workflow}，而非 Autonomous Agent。但这并不意味着 Agentic Workflow 一定比 Autonomous Agent 差——Agentic 是一个光谱，Autonomous Agent 是终极形态。
\end{knowledgebox}

预定义 Workflow 的优势在于\textbf{可控、可靠、可预测}。完全自主的 Agent（纯 ReAct 范式）受限于模型能力，可能表现不佳。

\subsection{OpenAI 关于 Reasoning Model 的最佳实践}

OpenAI API 官方文档中关于 Reasoning Model 的定位：
\begin{itemize}
    \item \textbf{O 系列推理模型}：适合做 Planner——可以对复杂任务执行更长时间、更深入的思考，更有效地制定策略
    \item \textbf{GPT 非推理模型}：适合做 Executor——追求低延迟、高成本效益，直接执行具体任务
\end{itemize}

\subsection{本章小结}

Agentic Workflow 胜在可控可靠，是当前实际应用中的主流选择。强模型做规划、弱模型做执行的策略与 OpenAI 的官方建议一致。

